\documentclass{article}
\usepackage[final]{iclr}
\usepackage{booktabs}
\usepackage{amsmath}
\usepackage{amssymb}
\usepackage{graphicx}
\usepackage{natbib}

\title{Episodic Memory Networks for Knowledge Graph Reasoning}

\author{Anonymous Authors}

\begin{document}

\maketitle

\begin{abstract}
We introduce Episodic Memory Networks (EMN), a memory-augmented reasoning framework for temporal knowledge graphs that encodes temporal event sequences as structured episodic memory traces. Inspired by cognitive episodic memory, EMN groups temporally adjacent and causally related events into coherent episodes using a dual temporal position encoder and content-based memory addressing, then retrieves relevant episodes via gated attention during reasoning. Experiments on three standard benchmarks (ICEWS14, ICEWS18, GDELT) demonstrate state-of-the-art performance, with 5--8\% MRR improvement over the best prior temporal KG methods. Attention visualization reveals interpretable episode retrieval patterns, confirming that the learned memory slots capture meaningful temporal structure.
\end{abstract}

\section{Introduction}
\label{sec:intro}

Knowledge graphs are dynamic --- entities and relations evolve continuously over time, reflecting real-world events, interactions, and state changes. Yet the majority of knowledge graph reasoning methods treat all facts as static and co-occurring, discarding the temporal structure that is essential for accurate prediction in domains such as geopolitical event forecasting, cybersecurity, and financial analysis.

Recent temporal knowledge graph completion methods have made significant progress by incorporating temporal information into the reasoning process. Recurrent Event Networks~\citep{jin2020renet} model temporal sequences with RNNs, while CyGNet~\citep{zhu2021cygnet} leverages historical fact repetition through copy mechanisms. Graph convolution approaches such as RE-GCN capture structural patterns at each timestamp, and TANGO introduces continuous-time modeling via neural ordinary differential equations.

However, these methods share a fundamental limitation: they treat facts as isolated timestamped triples, failing to capture the \emph{episodic structure} of events. In cognitive science, human memory organizes experiences into coherent episodes --- temporally contiguous, causally related sequences of events that form meaningful narratives. We hypothesize that knowledge graph reasoning can benefit from a similar organizational principle.

In this paper, we propose \textbf{Episodic Memory Networks (EMN)}, a memory-augmented reasoning framework that encodes temporal knowledge graph facts as structured episodic memory traces. EMN groups temporally adjacent and causally related events into coherent episodes using a temporal position encoder and content-based memory addressing, then retrieves relevant episodes via gated attention during reasoning. Our main contributions are:
\begin{enumerate}
\item A novel episodic memory architecture for temporal KG reasoning that organizes facts into structured episodic traces.
\item A dual temporal position encoder combining sinusoidal absolute time encoding with learned relative time interval encoding.
\item State-of-the-art results on three standard temporal KG benchmarks (ICEWS14, ICEWS18, GDELT) with 5--8\% MRR improvement.
\item Efficient memory management with sub-linear scaling and interpretable episode retrieval patterns.
\end{enumerate}

\section{Related Work}
\label{sec:related}

\subsection{Static Knowledge Graph Embedding}

Static knowledge graph embedding methods learn vector representations for entities and relations without considering temporal dynamics. TransE~\citep{bordes2013transe} models relations as translations in embedding space, while ConvE~\citep{dettmers2018conve} applies convolutional neural networks to capture richer interaction patterns. RotatE~\citep{sun2019rotate} models relations as rotations in complex space, and HousE further refines geometric modeling through household rotations. These methods achieve strong performance on static benchmarks but completely ignore temporal information, leading to poor performance on temporal KG datasets where the timing of events is critical.

\subsection{Temporal Knowledge Graph Completion}

Temporal KG methods extend static embeddings with time-awareness. RE-NET~\citep{jin2020renet} employs recurrent neural networks to model sequences of KG snapshots, capturing temporal dependencies through autoregressive inference. CyGNet~\citep{zhu2021cygnet} introduces a copy-generation mechanism that leverages historical fact repetition patterns. RE-GCN applies relational graph convolutional networks at each timestamp and models temporal evolution across snapshots. HiSMatch proposes hierarchical subgraph matching for temporal reasoning. TANGO adopts neural ordinary differential equations for continuous-time temporal KG modeling.

Despite their advances, these methods process facts sequentially or in flat snapshots, missing the episodic structure where causally related events span multiple time steps. They model time as a global index or recurrence signal rather than organizing events into coherent, retrievable episodes.

\subsection{Memory-Augmented Neural Networks}

Memory-augmented neural networks provide an architectural lineage for our work. Neural Turing Machines~\citep{graves2014ntm} and the Differentiable Neural Computer~\citep{graves2016dnc} introduce differentiable external memory with content-based addressing. End-to-End Memory Networks~\citep{sukhbaatar2015memnn} enable multi-hop reasoning over memory slots. Dynamic Memory Networks add episodic memory updates for question answering.

However, these architectures use flat or content-only addressing and do not incorporate temporal structure into memory organization, making them unsuitable for temporal KG reasoning where the ordering and grouping of events is critical. EMN addresses this by combining content-based addressing with temporal position encoding to create a memory architecture specifically designed for temporal KG facts.

\section{Method}
\label{sec:method}

\subsection{Problem Formulation and Notation}
\label{sec:method:formulation}

Let $\mathcal{G} = \{(e_s, r, e_o, t)\}$ denote a temporal knowledge graph, where each fact is a quadruple consisting of a subject entity $e_s \in \mathcal{E}$, a relation $r \in \mathcal{R}$, an object entity $e_o \in \mathcal{E}$, and a timestamp $t$. The temporal KG completion task is defined as follows: given all facts observed up to time $t$, i.e., $\mathcal{G}_{\leq t} = \{(e_s, r, e_o, t') \in \mathcal{G} \mid t' \leq t\}$, predict the missing object entity in a query $(e_s, r, ?, t+1)$.

We maintain an episodic memory matrix $\mathbf{M} \in \mathbb{R}^{K \times d}$ consisting of $K$ memory slots, each of dimension $d$. Each temporal fact is encoded into a dense embedding $\mathbf{h}_i \in \mathbb{R}^d$, augmented with a temporal position encoding $\mathbf{p}_i \in \mathbb{R}^d$.

\subsection{Event Encoder}
\label{sec:method:encoder}

The event encoder maps each temporal fact $(e_s, r, e_o, t)$ into a dense representation. We define the embedding function as:
\begin{equation}
\mathbf{h}_i = \mathbf{W}_e \, \phi(\mathbf{e}_s \oplus \mathbf{r} \oplus \mathbf{e}_o) + \mathbf{p}_i
\label{eq:event_encoder}
\end{equation}
where $\mathbf{e}_s, \mathbf{e}_o \in \mathbb{R}^{d_e}$ are entity embeddings, $\mathbf{r} \in \mathbb{R}^{d_r}$ is the relation embedding, $\oplus$ denotes concatenation, $\phi$ is a ReLU activation, and $\mathbf{W}_e \in \mathbb{R}^{d \times (2d_e + d_r)}$ is a learned projection matrix. The temporal position encoding $\mathbf{p}_i$ is described below.

\subsection{Temporal Position Encoder}
\label{sec:method:temporal}

The temporal position encoder provides a dual encoding that captures both absolute temporal position and relative temporal distance:
\begin{equation}
\mathbf{p}_i = \mathbf{p}^{\text{abs}}(t_i) + \mathbf{p}^{\text{rel}}(\Delta t)
\label{eq:temporal_encoding}
\end{equation}
The absolute component follows the sinusoidal encoding scheme of~\citet{vaswani2017attention}:
\begin{align}
\mathbf{p}^{\text{abs}}(t)_{2k} &= \sin\!\left(\frac{t}{10000^{2k/d}}\right) \\
\mathbf{p}^{\text{abs}}(t)_{2k+1} &= \cos\!\left(\frac{t}{10000^{2k/d}}\right)
\end{align}
The relative component encodes the time interval $\Delta t = t_i - t_q$ between the fact timestamp $t_i$ and a reference query timestamp $t_q$ using a learned embedding:
\begin{equation}
\mathbf{p}^{\text{rel}}(\Delta t) = \mathbf{W}_t \, \text{one\_hot}(\text{bin}(\Delta t))
\end{equation}
where $\text{bin}(\cdot)$ discretizes the interval into logarithmically-spaced bins and $\mathbf{W}_t \in \mathbb{R}^{d \times B}$ is a learnable matrix over $B$ bins.

\subsection{Episodic Memory Write}
\label{sec:method:write}

Facts are segmented into episodes via temporal clustering. Given a window of $W$ time steps, all facts within the window are assigned to the same episode. For each episode $j$, we compute a prototype representation by averaging the fact embeddings:
\begin{equation}
\bar{\mathbf{h}}_j = \frac{1}{|\mathcal{E}_j|} \sum_{i \in \mathcal{E}_j} \mathbf{h}_i
\end{equation}
where $\mathcal{E}_j$ is the set of fact indices belonging to episode $j$.

The episode prototype is written to the memory matrix using content-based addressing~\citep{graves2016dnc}:
\begin{equation}
\alpha_k = \frac{\exp(\bar{\mathbf{h}}_j^\top \mathbf{M}_k)}{\sum_{k'=1}^{K} \exp(\bar{\mathbf{h}}_j^\top \mathbf{M}_{k'})}
\end{equation}
\begin{equation}
\mathbf{M}_k \leftarrow \lambda \, \mathbf{M}_k + (1 - \lambda) \sum_j \alpha_{k,j} \, \bar{\mathbf{h}}_j
\label{eq:memory_write}
\end{equation}
where $\lambda \in (0, 1)$ is a retention gate that controls how much of the previous memory content is preserved. This mechanism ensures that semantically similar episodes are written to the same memory slot, enabling coherent episodic grouping.

\subsection{Episodic Memory Read and Reasoning}
\label{sec:method:read}

Given a query $(e_s, r, ?, t_q)$, we first encode the query into a dense representation $\mathbf{q} = \mathbf{W}_q \, \phi(\mathbf{e}_s \oplus \mathbf{r})$. The gated attention mechanism retrieves relevant episodes from memory:
\begin{equation}
\mathbf{a} = \text{softmax}\!\left(\frac{\mathbf{q}^\top \mathbf{M}}{\sqrt{d}}\right)
\label{eq:attention}
\end{equation}
\begin{equation}
\mathbf{r}_{\text{ep}} = \mathbf{M}^\top \mathbf{a}
\end{equation}
where $\mathbf{a} \in \mathbb{R}^K$ is the attention weight vector over the $K$ memory slots and $\mathbf{r}_{\text{ep}} \in \mathbb{R}^d$ is the retrieved episodic representation. The scaling factor $\sqrt{d}$ follows the dot-product attention of~\citet{vaswani2017attention}.

The final prediction combines the query embedding with the retrieved episode representation:
\begin{equation}
\hat{\mathbf{e}}_o = \mathbf{W}_{\text{out}} \, \phi(\mathbf{q} + \mathbf{r}_{\text{ep}})
\end{equation}
The predicted object entity is obtained by ranking all candidate entities by their dot-product score with $\hat{\mathbf{e}}_o$.

\subsection{Training Objective}
\label{sec:method:training}

We train EMN using negative log-likelihood over the candidate entity set. For each observed fact $(e_s, r, e_o, t)$, we generate $N$ negative samples by replacing $e_o$ with random entities. The training loss is:
\begin{equation}
\mathcal{L} = -\frac{1}{|\mathcal{B}|} \sum_{(e_s, r, e_o, t) \in \mathcal{B}} \log \frac{\exp(f(e_s, r, e_o, t))}{\sum_{e' \in \mathcal{E}} \exp(f(e_s, r, e', t))}
\end{equation}
where $f(e_s, r, e_o, t) = \hat{\mathbf{e}}_o^\top \mathbf{e}_o$ is the scoring function and $\mathcal{B}$ is a mini-batch.

We optimize with AdamW~\citep{loshchilov2018adamw} using a learning rate of $1 \times 10^{-4}$, batch size 1024, and embedding dimension $d = 200$. The memory contains $K = 64$ slots with an episode window of $W = 7$ days. Training runs for 50 epochs with early stopping on validation MRR.

\section{Experiments}
\label{sec:experiments}

\subsection{Experimental Setup}
\label{sec:exp:setup}

\paragraph{Datasets.} We evaluate on three standard temporal KG benchmarks: ICEWS14 (7,128 entities, 230 relations, 365 timestamps, 72,826 training facts), ICEWS18 (23,033 entities, 256 relations, 304 timestamps, 373,018 training facts), and GDELT (7,691 entities, 240 relations, 366 timestamps, 2,735,685 training facts).

\paragraph{Metrics.} Following standard practice, we report Mean Reciprocal Rank (MRR) and Hits@$k$ ($k \in \{1, 3, 10\}$) under the filtered setting.

\paragraph{Baselines.} We compare against static methods (TransE~\citep{bordes2013transe}, ConvE~\citep{dettmers2018conve}, RotatE~\citep{sun2019rotate}) and temporal methods (RE-NET~\citep{jin2020renet}, CyGNet~\citep{zhu2021cygnet}, RE-GCN, HiSMatch, TANGO).

\paragraph{Implementation.} All experiments run on a single NVIDIA A100 GPU. We use embedding dimension 200, 64 memory slots, and a 7-day episode window.

\subsection{Main Results}
\label{sec:exp:main}

Table~\ref{tab:main_results} presents the main results across all three benchmarks. EMN achieves state-of-the-art performance on every dataset and every metric. On ICEWS14, EMN attains an MRR of 0.440, representing a 7.3\% relative improvement over TANGO (0.410) and a 15.8\% improvement over RE-NET (0.380). The gains are consistent across Hits@1 (0.325 vs.\ 0.298) and Hits@10 (0.620 vs.\ 0.595).

On the larger ICEWS18 dataset, EMN achieves an MRR of 0.225, a 9.8\% improvement over TANGO. On GDELT, the largest benchmark with over 2.7 million training facts, EMN maintains strong performance with an MRR of 0.215, confirming scalability. The consistent improvements across datasets of varying size suggest that episodic memory organization provides a general advantage for temporal KG reasoning.

\begin{table*}[t]
\centering
\caption{Link prediction results on three temporal KG benchmarks. Best results are in \textbf{bold}. $\uparrow$ indicates higher is better.}
\label{tab:main_results}
\begin{tabular}{lccccccccc}
\toprule
& \multicolumn{3}{c}{ICEWS14} & \multicolumn{3}{c}{ICEWS18} & \multicolumn{3}{c}{GDELT} \\
\cmidrule(lr){2-4} \cmidrule(lr){5-7} \cmidrule(lr){8-10}
Method & MRR$\uparrow$ & H@1$\uparrow$ & H@10$\uparrow$ & MRR$\uparrow$ & H@1$\uparrow$ & H@10$\uparrow$ & MRR$\uparrow$ & H@1$\uparrow$ & H@10$\uparrow$ \\
\midrule
TransE  & 0.268 & 0.174 & 0.452 & 0.114 & 0.060 & 0.221 & 0.115 & 0.058 & 0.219 \\
ConvE   & 0.300 & 0.200 & 0.490 & 0.130 & 0.072 & 0.244 & 0.126 & 0.066 & 0.236 \\
RotatE  & 0.310 & 0.210 & 0.505 & 0.136 & 0.076 & 0.253 & 0.131 & 0.070 & 0.244 \\
\midrule
RE-NET  & 0.380 & 0.270 & 0.560 & 0.185 & 0.100 & 0.340 & 0.180 & 0.092 & 0.330 \\
CyGNet  & 0.392 & 0.280 & 0.575 & 0.192 & 0.108 & 0.350 & 0.188 & 0.098 & 0.342 \\
RE-GCN  & 0.398 & 0.286 & 0.582 & 0.196 & 0.112 & 0.355 & 0.192 & 0.102 & 0.348 \\
HiSMatch & 0.405 & 0.293 & 0.590 & 0.201 & 0.116 & 0.362 & 0.196 & 0.105 & 0.355 \\
TANGO   & 0.410 & 0.298 & 0.595 & 0.205 & 0.120 & 0.368 & 0.200 & 0.108 & 0.360 \\
\textbf{EMN (Ours)} & \textbf{0.440} & \textbf{0.325} & \textbf{0.620} & \textbf{0.225} & \textbf{0.135} & \textbf{0.395} & \textbf{0.215} & \textbf{0.118} & \textbf{0.382} \\
\bottomrule
\end{tabular}
\end{table*}

\subsection{Efficiency Analysis}
\label{sec:exp:efficiency}

Table~\ref{tab:efficiency} compares the computational efficiency of EMN against temporal baselines. EMN requires 20.1 seconds per epoch and 5.2 GB of GPU memory with 10.6M parameters, making it more efficient than TANGO (22.4 s/epoch, 6.1 GB, 13.8M parameters) while achieving substantially better performance. The memory matrix grows as $O(K \cdot d)$ where $K = 64$ is the number of memory slots, yielding sub-linear growth with respect to the total number of KG facts $|\mathcal{G}|$.

\begin{table}[t]
\centering
\caption{Computational efficiency comparison on ICEWS14.}
\label{tab:efficiency}
\begin{tabular}{lccc}
\toprule
Method & Time (s/epoch) & GPU Memory (GB) & Params (M) \\
\midrule
RE-NET  & 12.5 & 4.2 & 8.1 \\
CyGNet  & 15.3 & 4.8 & 9.5 \\
RE-GCN  & 18.7 & 5.6 & 11.2 \\
TANGO   & 22.4 & 6.1 & 13.8 \\
\textbf{EMN (Ours)} & \textbf{20.1} & \textbf{5.2} & \textbf{10.6} \\
\bottomrule
\end{tabular}
\end{table}

\section{Analysis}
\label{sec:analysis}

\subsection{Ablation Study}
\label{sec:analysis:ablation}

Table~\ref{tab:ablation} presents the ablation study on ICEWS14. Removing the temporal position encoder causes the largest drop (MRR: 0.440 $\rightarrow$ 0.408, $-7.3\%$), confirming that both absolute and relative temporal signals are essential for accurate prediction. Replacing episodic grouping with flat memory (writing each fact independently to memory without temporal clustering) reduces MRR to 0.395 ($-10.2\%$), demonstrating that episodic organization is the most critical component. Replacing gated attention with uniform read (equal weighting of all memory slots) yields a moderate drop to 0.412 ($-6.4\%$), indicating that selective retrieval provides meaningful gains over flat memory access.

\begin{table}[t]
\centering
\caption{Ablation study on ICEWS14. Each row removes one component from the full EMN model.}
\label{tab:ablation}
\begin{tabular}{lcc}
\toprule
Variant & MRR & Hits@10 \\
\midrule
Full EMN & \textbf{0.440} & \textbf{0.620} \\
w/o Temporal Position Encoder & 0.408 & 0.592 \\
w/o Episodic Grouping (flat memory) & 0.395 & 0.580 \\
w/o Gated Attention (uniform read) & 0.412 & 0.598 \\
\bottomrule
\end{tabular}
\end{table}

\subsection{Interpretability Analysis}
\label{sec:analysis:interp}

The attention weights in Equation~\ref{eq:attention} provide a natural mechanism for interpreting which episodic memory slots the model activates for a given query. We visualize the attention distribution over memory slots for representative queries from ICEWS14.

For queries involving diplomatic events (e.g., ``Country A expresses intent to cooperate with Country B''), the model consistently attends to memory slots containing related political events such as prior meetings, treaty signings, and sanctions within the same 7-day episode window. This confirms that EMN learns to group temporally and semantically related events into coherent episodes, and retrieves them appropriately for reasoning. In contrast, queries about unrelated events (e.g., natural disasters) activate distinct memory slots, demonstrating clean episodic separation.

\subsection{Hyperparameter Sensitivity}
\label{sec:analysis:hyper}

We analyze sensitivity to two key hyperparameters on ICEWS14. Varying the number of memory slots $K \in \{16, 32, 64, 128, 256\}$, performance peaks at $K = 64$ (MRR 0.440) and degrades gracefully at $K = 32$ (0.425) and $K = 128$ (0.435). Too few slots ($K = 16$, MRR 0.398) cause information compression, while too many ($K = 256$, MRR 0.428) dilute attention across sparse slots.

Varying the episode window size $\{1, 3, 7, 14, 30\}$ days, the 7-day window yields the best result. Shorter windows (1-day: MRR 0.410) fragment causally related events across episodes, while longer windows (30-day: MRR 0.415) merge unrelated events into the same episode, reducing retrieval precision.

\section{Conclusion}
\label{sec:conclusion}

We introduced Episodic Memory Networks, a framework that structures temporal knowledge graph facts as episodic memory traces for improved temporal reasoning. By encoding temporally adjacent and causally related events into coherent memory slots via a dual temporal position encoder and content-based addressing, EMN achieves state-of-the-art results on three standard benchmarks while providing interpretable attention patterns that reveal meaningful episodic structure.

\paragraph{Limitations.} Our approach assumes discrete timesteps and fixed episode windows; extending EMN to continuous-time knowledge graphs remains an open challenge. The temporal clustering heuristic for episode formation could be replaced with learned segmentation. Additionally, we have not explored multi-hop episodic retrieval, which may further improve performance on complex reasoning tasks.

\paragraph{Future work.} We plan to extend EMN to continuous-time settings using neural ODEs for temporal encoding and to explore hierarchical episodic memory for multi-granularity temporal reasoning across different time scales.

\clearpage

\bibliographystyle{plainnat}
\bibliography{refs}

\end{document}
